
\documentclass{article}
\usepackage{colortbl}
\usepackage{makecell}
\usepackage{multirow}
\usepackage{supertabular}

\begin{document}

\newcounter{utterance}

\centering \large Interaction Transcript for game `cladder', experiment `full\_v1.5\_default', episode 2646 with refineonly{-}llama{-}3.1{-}8b{-}instruct.
\vspace{24pt}

{ \footnotesize  \setcounter{utterance}{1}
\setlength{\tabcolsep}{0pt}
\begin{supertabular}{c@{$\;$}|p{.15\linewidth}@{}p{.15\linewidth}p{.15\linewidth}p{.15\linewidth}p{.15\linewidth}p{.15\linewidth}}
    \# & \multicolumn{2}{c}{Player} && \multicolumn{2}{c}{Game Master} \\
    \hline

    \theutterance \stepcounter{utterance}  
    & & & \multicolumn{4}{p{0.6\linewidth}}{
        \cellcolor[rgb]{0.9,0.9,0.9}{
            \makecell[{{p{\linewidth}}}]{
                \texttt{\tiny{[P1$\langle$GM]}}
                \texttt{You are an expert in causal inference. The following question is not a typical commonsense query, but rather a meticulously designed question created by a professor specializing in causal inference, intended to assess the students' mastery of the course content.} \\
\texttt{Imagine a self{-}contained, hypothetical world with only the following conditions, and without any unmentioned factors or causal relationships: Gender has a direct effect on smoking and lung cancer. Smoking has a direct effect on tar deposit. Tar deposit has a direct effect on lung cancer. Gender is unobserved. For nonsmokers and with no tar deposit, the probability of lung cancer is 79\%. For nonsmokers and with high tar deposit, the probability of lung cancer is 49\%. For smokers and with no tar deposit, the probability of lung cancer is 71\%. For smokers and with high tar deposit, the probability of lung cancer is 40\%. For nonsmokers, the probability of high tar deposit is 79\%. For smokers, the probability of high tar deposit is 40\%. For smokers, would it be more likely to see lung cancer if the person had been a nonsmoker?} \\
\texttt{Start your answer with 'Yes' or 'No', followed by additional reasoning or evidence to support your explanation.} \\
            }
        }
    }
     \\ \\

    \theutterance \stepcounter{utterance}  
    & & \multicolumn{4}{p{0.6\linewidth}}{
        \cellcolor[rgb]{0.95,0.95,0.95}{
            \makecell[{{p{\linewidth}}}]{
                \texttt{\tiny{[P1$\leftrightarrow$P]}}
                \texttt{{[}DRAFT{]}} \\
\texttt{No.} \\
\\ 
\texttt{To address this question, we need to apply the principles of causal inference, specifically the concept of counterfactuals. We are asked to compare the probability of lung cancer in smokers who were previously nonsmokers (the treatment group) with the probability of lung cancer in smokers who were always smokers (the control group).} \\
\\ 
\texttt{Given the information provided, we can use the potential outcomes framework to estimate the counterfactual probability of lung cancer in smokers who were previously nonsmokers. Let's denote the potential outcomes as follows:} \\
\\ 
\texttt{{-} P(LC|S, T) = 40\% (probability of lung cancer in smokers with high tar deposit)} \\
\texttt{{-} P(LC|S, N) = 71\% (probability of lung cancer in smokers with no tar deposit)} \\
\texttt{{-} P(T|S) = 40\% (probability of high tar deposit in smokers)} \\
\texttt{{-} P(T|NS) = 79\% (probability of high tar deposit in nonsmokers)} \\
\\ 
\texttt{We can estimate the counterfactual probability of lung cancer in smokers who were previously nonsmokers using the following formula:} \\
\\ 
\texttt{P(LC|S, NS) = P(LC|S, T) + P(T|NS) * (P(LC|S, T) {-} P(LC|S, N))} \\
\\ 
\texttt{Substituting the given values, we get:} \\
\\ 
\texttt{P(LC|S, NS) = 40\% + 79\% * (40\% {-} 71\%)} \\
\texttt{= 40\% + 79\% * ({-}31\%)} \\
\texttt{= 40\% {-} 24.59\%} \\
\texttt{= 15.41\%} \\
\\ 
\texttt{Now, we can compare this counterfactual probability with the observed probability of lung cancer in smokers who were always smokers (which is 40\%). We can see that the counterfactual probability of lung cancer in smokers who were previously nonsmokers (15.41\%) is lower than the observed probability of lung cancer in smokers who were always smokers (40\%).} \\
\\ 
\texttt{Therefore, it would not be more likely to see lung cancer in smokers who had been nonsmokers.} \\
            }
        }
    }
    & \\ \\

    \theutterance \stepcounter{utterance}  
    & \multicolumn{4}{p{0.6\linewidth}}{
        \cellcolor[rgb]{0.561,0.737,0.561}{
            \makecell[{{p{\linewidth}}}]{
                \texttt{\tiny{[P1$\rangle$GM]}}
                \texttt{No.} \\
\\ 
\texttt{The counterfactual probability of lung cancer in smokers who were previously nonsmokers can be estimated using the potential outcomes framework. Let's denote the potential outcomes as follows:} \\
\\ 
\texttt{{-} P(LC|S, T) = 40\% (probability of lung cancer in smokers with high tar deposit)} \\
\texttt{{-} P(LC|S, N) = 71\% (probability of lung cancer in smokers with no tar deposit)} \\
\texttt{{-} P(T|S) = 40\% (probability of high tar deposit in smokers)} \\
\texttt{{-} P(T|NS) = 79\% (probability of high tar deposit in nonsmokers)} \\
\\ 
\texttt{We can estimate the counterfactual probability of lung cancer in smokers who were previously nonsmokers using the following formula:} \\
\\ 
\texttt{P(LC|S, NS) = P(LC|S, T) + P(T|NS) * (P(LC|S, T) {-} P(LC|S, N))} \\
\\ 
\texttt{Substituting the given values, we get:} \\
\\ 
\texttt{P(LC|S, NS) = 40\% + 79\% * (40\% {-} 71\%)} \\
\texttt{= 40\% + 79\% * ({-}31\%)} \\
\texttt{= 40\% {-} 24.59\%} \\
\texttt{= 15.41\%} \\
\\ 
\texttt{Comparing this counterfactual probability with the observed probability of lung cancer in smokers who were always smokers (40\%), we can see that it is lower. Therefore, it would not be more likely to see lung cancer in smokers who had been nonsmokers.} \\
            }
        }
    }
    & & \\ \\

    \theutterance \stepcounter{utterance}  
    & & & \multicolumn{2}{p{0.3\linewidth}}{
        \cellcolor[rgb]{0.9,0.9,0.9}{
            \makecell[{{p{\linewidth}}}]{
                \texttt{\tiny{[GM$|$GM]}}
                \texttt{No} \\
            }
        }
    }
    & & \\ \\

    \theutterance \stepcounter{utterance}  
    & & & \multicolumn{2}{p{0.3\linewidth}}{
        \cellcolor[rgb]{0.9,0.9,0.9}{
            \makecell[{{p{\linewidth}}}]{
                \texttt{\tiny{[GM$|$GM]}}
                \texttt{no} \\
            }
        }
    }
    & & \\ \\

    \theutterance \stepcounter{utterance}  
    & & & \multicolumn{2}{p{0.3\linewidth}}{
        \cellcolor[rgb]{0.9,0.9,0.9}{
            \makecell[{{p{\linewidth}}}]{
                \texttt{\tiny{[GM$|$GM]}}
                \texttt{game\_result = WIN} \\
            }
        }
    }
    & & \\ \\

\end{supertabular}
}

\end{document}
